% =========================================================================
% Capítulo 4: Cuerpo rígido
% Tema 01: Cuerpo rígido
% Problema 21: Disco rodante sobre una cuerda con resorte y amortiguador cúbico
% Origen: Ayudantía 12 (Sólido Rígido), 2026; corresponde al Ejemplo 5-5 del texto de referencia
% =========================================================================
\subsubsection{Disco rodante sobre una cuerda con resorte y amortiguador cúbico}

Un disco homogéneo de masa $m$ y radio $r$ rueda sin deslizar en el plano vertical a lo largo de una cuerda inextensible, como se muestra en la Figura \ref{fig:c4_e01_p21_disco_cuerda}. El centro del disco está ubicado en el punto $C$, mientras que los puntos $A$ y $B$ se encuentran en los extremos del diámetro horizontal del disco. La cuerda está unida en un extremo a un punto fijo $O$ y en su otro extremo a un resorte lineal con constante elástica $K$ y a un amortiguador con constante de amortiguamiento $c$. El amortiguador ejerce una fuerza en dirección opuesta a la velocidad en el punto $B$ y es proporcional al cubo de la magnitud de la velocidad en dicho punto, es decir, la fuerza del amortiguador tiene la forma:
\[
\vec{F}_d = -c \|\vec{v}_B\|^3 \frac{\vec{v}_B}{\|\vec{v}_B\|}
\]
Asumiendo que la cuerda permanece siempre tensa, que el resorte no está estirado con una longitud natural $L$ cuando el ángulo $\theta$ es cero, y que la gravedad actúa hacia abajo:

\begin{enumerate}
    \item Determine la ecuación diferencial de movimiento para el disco en términos del ángulo $\theta$.
\end{enumerate}

\begin{figure}[!ht]
    \centering
    \begin{tikzpicture}[>=Latex, scale=1.1]
        % Soportes superiores (Techos O y P)
        \fill[pattern=north east lines] (-2.2, 4.5) rectangle (-0.8, 4.75);
        \draw[thick] (-2.2, 4.5) -- (-0.8, 4.5);
        \node[below left] at (-1.5, 4.5) {$O$};

        \fill[pattern=north east lines] (0.8, 4.5) rectangle (2.2, 4.75);
        \draw[thick] (0.8, 4.5) -- (2.2, 4.5);
        \node[below right] at (1.5, 4.5) {$P$};

        % Disco homogéneo con sombreado radial
        \shadedraw[inner color=white, outer color=gray!45, draw=black, thick] (0,0) circle (1.5cm);
    
        % Centro C y radio r
        \fill (0,0) circle (2pt) node[left] {$C$};
        \draw[->, thick] (0,0) -- (45:1.5) node[midway, below right] {$r$};

        % Puntos A y B
        \node[right] at (-1.5, 0) {$A$};
        \node[left] at (1.5, 0) {$B$};

        % Cuerda izquierda (O a A)
        \draw[thick] (-1.5, 4.5) -- (-1.5, 0);

        % Conexión derecha (B al sistema resorte-amortiguador y a P)
        \draw[thick] (1.5, 0) -- (1.5, 1.3);
        \draw[thick] (1.2, 1.3) -- (1.8, 1.3); % Barra inferior
        \draw[thick] (1.2, 3.2) -- (1.8, 3.2); % Barra superior
        \draw[thick] (1.5, 3.2) -- (1.5, 4.5);

        % Amortiguador (izquierda del bloque paralelo)
        \draw[thick] (1.32, 1.3) -- (1.32, 2.15);
        \draw[thick] (1.22, 2.15) -- (1.42, 2.15); % Pistón
        \draw[thick] (1.20, 2.0) -- (1.20, 2.35) -- (1.44, 2.35) -- (1.44, 2.0); % Cilindro
        \draw[thick] (1.32, 2.35) -- (1.32, 3.2);
        \node[left] at (1.18, 2.25) {$c$};

        % Resorte (derecha del bloque paralelo)
        \draw[thick, decorate, decoration={coil, aspect=0.4, segment length=3.2pt, amplitude=3.2pt}] (1.68, 1.3) -- (1.68, 3.2);
        \node[right] at (1.82, 2.25) {$K$};

        % Flecha de rotación theta
        \draw[->, thick] (285:1.85) arc (285:255:1.85) node[midway, below] {$\theta$};

        % Vector gravedad
        \draw[->, thick] (2.8, 4.2) -- (2.8, 3.2) node[midway, left] {$\vec{g}$};
    \end{tikzpicture}

    \caption{Disco que rueda sin deslizar sobre una cuerda unida a un resorte y un amortiguador.}
    \label{fig:c4_e01_p21_disco_cuerda}
\end{figure}

% --- SOLUCIÓN ---------------------------------------------------------------
\begin{solucion}{Disco rodante sobre una cuerda con resorte y amortiguador cúbico}
\subsubsection*{Cinemática}

Elegimos un marco de referencia inercial fijo $\mathcal{F}$ con el siguiente sistema de coordenadas:
\begin{itemize}
    \item Origen en $C$ cuando $\theta = 0$
    \item $\vec{E}_x =$ Verticalmente hacia abajo
    \item $\vec{E}_z =$ Hacia adentro de la página
    \item $\vec{E}_y = \vec{E}_z \times \vec{E}_x$ (Horizontalmente hacia la izquierda)
\end{itemize}

Como el disco rueda sin deslizar sobre la cuerda y el punto $O$ es fijo, cada punto de la cuerda entre $O$ y $A$ tiene velocidad cero en $\mathcal{F}$. En particular:
\begin{equation}
    {}^{\mathcal{F}}\vec{v}_A = \vec{0}
\end{equation}
La velocidad y aceleración angular del disco $\mathcal{R}$ en el marco $\mathcal{F}$ son:
\begin{equation}
    {}^{\mathcal{F}}\vec{\omega}^{\mathcal{R}} = \dot{\theta}\vec{E}_z, \qquad {}^{\mathcal{F}}\vec{\alpha}^{\mathcal{R}} = \ddot{\theta}\vec{E}_z
\end{equation}
Por geometría, las posiciones relativas entre $A$, $B$ y $C$ son:
\begin{equation}
    \vec{r}_C - \vec{r}_A = -r\vec{E}_y \implies \vec{r}_A - \vec{r}_C = r\vec{E}_y, \qquad \vec{r}_B - \vec{r}_A = -2r\vec{E}_y
\end{equation}
Calculamos las velocidades de $C$ y $B$ a partir de ${}^{\mathcal{F}}\vec{v}_A = \vec{0}$:
\begin{align}
    {}^{\mathcal{F}}\vec{v}_C &= {}^{\mathcal{F}}\vec{v}_A + {}^{\mathcal{F}}\vec{\omega}^{\mathcal{R}} \times (\vec{r}_C - \vec{r}_A) = \dot{\theta}\vec{E}_z \times (-r\vec{E}_y) = r\dot{\theta}\vec{E}_x \\
    {}^{\mathcal{F}}\vec{v}_B &= {}^{\mathcal{F}}\vec{v}_A + {}^{\mathcal{F}}\vec{\omega}^{\mathcal{R}} \times (\vec{r}_B - \vec{r}_A) = \dot{\theta}\vec{E}_z \times (-2r\vec{E}_y) = 2r\dot{\theta}\vec{E}_x
\end{align}
Derivando ${}^{\mathcal{F}}\vec{v}_C$, obtenemos la aceleración del centro de masa $C$:
\begin{equation}
    {}^{\mathcal{F}}\vec{a}_C = \frac{{}^{\mathcal{F}}d}{dt}({}^{\mathcal{F}}\vec{v}_C) = r\ddot{\theta}\vec{E}_x
\end{equation}
Calculamos la aceleración del punto instantáneo de contacto $A$ en el disco:
\begin{align}
    {}^{\mathcal{F}}\vec{a}_A &= {}^{\mathcal{F}}\vec{a}_C + {}^{\mathcal{F}}\vec{\alpha}^{\mathcal{R}} \times (\vec{r}_A - \vec{r}_C) + {}^{\mathcal{F}}\vec{\omega}^{\mathcal{R}} \times [{}^{\mathcal{F}}\vec{\omega}^{\mathcal{R}} \times (\vec{r}_A - \vec{r}_C)] \nonumber \\
    &= r\ddot{\theta}\vec{E}_x + \ddot{\theta}\vec{E}_z \times r\vec{E}_y + \dot{\theta}\vec{E}_z \times (\dot{\theta}\vec{E}_z \times r\vec{E}_y) \nonumber \\
    &= r\ddot{\theta}\vec{E}_x - r\ddot{\theta}\vec{E}_x - r\dot{\theta}^2\vec{E}_y = -r\dot{\theta}^2\vec{E}_y
\end{align}
Notamos que aunque ${}^{\mathcal{F}}\vec{v}_A = \vec{0}$, su aceleración no es nula (${}^{\mathcal{F}}\vec{a}_A \neq \vec{0}$), por lo que $A$ no es un punto fijo inercialmente.

\subsubsection*{Cinética}

Elegimos el punto $A$ como punto de referencia para eliminar la tensión desconocida $\vec{T}$ de la cuerda izquierda, aplicando la forma general de la 2ª Ley de Euler respecto a un punto arbitrario:
\begin{equation}
    \vec{M}_A - (\vec{r}_C - \vec{r}_A) \times m\,{}^{\mathcal{F}}\vec{a}_A = \frac{{}^{\mathcal{F}}d}{dt}({}^{\mathcal{F}}\vec{H}_A) \label{eq:c4_e01_p21_euler}
\end{equation}
El DCL del disco se muestra a continuación:

\begin{center}
\begin{tikzpicture}[scale=1.0, >=Stealth]
    \def\R{1.3}
    \shadedraw[ball color=gray!15, shading=ball, opacity=0.85, thick] (0,0) circle (\R cm);
    \fill (0,0) circle (2pt);
    % Peso
    \draw[->, very thick] (0,0) -- (0,-1.0) node[midway, left] {$m\vec{g}$};
    % Tensión en A
    \draw[->, very thick] (-\R, 0) -- (-\R, 1.8) node[left] {$\vec{T}$};
    % Fuerza resorte y amortiguador en B
    \draw[->, very thick] (\R, 0) -- (\R, 1.8) node[right] {$\vec{F}_s, \vec{F}_d$};
\end{tikzpicture}
\end{center}

Integrando ${}^{\mathcal{F}}\vec{v}_C = r\dot{\theta}\vec{E}_x$ y ${}^{\mathcal{F}}\vec{v}_B = 2r\dot{\theta}\vec{E}_x$ desde $\theta = 0$:
\begin{align}
    \vec{r}_C &= r\theta\vec{E}_x \\
    \vec{r}_A &= r\theta\vec{E}_x + r\vec{E}_y \\
    \vec{r}_B &= 2r\theta\vec{E}_x - r\vec{E}_y
\end{align}
El estiramiento del resorte es $l - l_0 = (2r\theta + L) - L = 2r\theta$ en la dirección $\vec{u}_s = \vec{E}_x$. Por tanto, las fuerzas que hacen momento respecto a $A$ son:
\begin{align}
    \vec{F}_s &= -K(l - l_0)\vec{u}_s = -2Kr\theta\vec{E}_x \\
    \vec{F}_d &= -c\|2r\dot{\theta}\vec{E}_x\|^3 \frac{2r\dot{\theta}\vec{E}_x}{\|2r\dot{\theta}\vec{E}_x\|} = -8cr^3\dot{\theta}^3\vec{E}_x \\
    m\vec{g} &= mg\vec{E}_x
\end{align}
Con los brazos de momento $\vec{r}_B - \vec{r}_A = -2r\vec{E}_y$ (o incluyendo su desplazamiento en $\vec{E}_x$, cuyo producto cruz con $\vec{E}_x$ es nulo) y $\vec{r}_C - \vec{r}_A = -r\vec{E}_y$, el momento neto respecto a $A$ es:
\begin{align}
    \vec{M}_A &= (-2r\vec{E}_y) \times (-2Kr\theta\vec{E}_x - 8cr^3\dot{\theta}^3\vec{E}_x) + (-r\vec{E}_y) \times (mg\vec{E}_x) \nonumber \\
    &= (-4Kr^2\theta - 16cr^4\dot{\theta}^3 + mgr)\vec{E}_z \quad (\text{usando } \vec{E}_y \times \vec{E}_x = -\vec{E}_z)
\end{align}

Calculamos el momento inercial del punto de referencia $A$:
\begin{equation}
    -(\vec{r}_C - \vec{r}_A) \times m\,{}^{\mathcal{F}}\vec{a}_A = -(-r\vec{E}_y) \times m(-r\dot{\theta}^2\vec{E}_y) = \vec{0}
\end{equation}

Finalmente, calculamos el momento angular respecto a $A$ con $\bar{I}_{zz} = \frac{1}{2}mr^2$:
\begin{align}
    {}^{\mathcal{F}}\vec{H}_A &= {}^{\mathcal{F}}\vec{H}_C + (\vec{r}_A - \vec{r}_C) \times m({}^{\mathcal{F}}\vec{v}_A - {}^{\mathcal{F}}\vec{v}_C) \nonumber \\
    &= \frac{1}{2}mr^2\dot{\theta}\vec{E}_z + r\vec{E}_y \times m(-r\dot{\theta}\vec{E}_x) = \frac{3}{2}mr^2\dot{\theta}\vec{E}_z
\end{align}
Derivando respecto al tiempo en $\mathcal{F}$:
\begin{equation}
    \frac{{}^{\mathcal{F}}d}{dt}({}^{\mathcal{F}}\vec{H}_A) = \frac{3}{2}mr^2\ddot{\theta}\vec{E}_z
\end{equation}

Sustituyendo en la ecuación (\ref{eq:c4_e01_p21_euler}) e igualando las componentes en $\vec{E}_z$:
\begin{equation}
    -4Kr^2\theta - 16cr^4\dot{\theta}^3 + mgr = \frac{3}{2}mr^2\ddot{\theta}
\end{equation}
Reordenando los términos, obtenemos la ecuación diferencial de movimiento del disco:
\respuesta{$\displaystyle \frac{3}{2}mr^2\ddot{\theta} + 16cr^4\dot{\theta}^3 + 4Kr^2\theta - mgr = 0$}
\end{solucion}
